% GENERATED FILE. Edit canonical lessons, then run npm run generate:books.
% canonical-source-hash: fnv1a64:ec148c623077b28d
% canonical-lessons: ML-C259-pavaykka, ML-C259-ventaykka, ML-C259-murinnakka, ML-C259-colam, ML-C259-katala

\chapter{Bitter gourd, Okra, Drumstick, Maize, Chickpeas}
\label{ch:ml-bitter-gourd-okra-drumstick-maize-chickpeas}
\hlchaptermodality{259}

\begin{chapteropening}
\textbf{By the end of this chapter:} \emph{I can say a bitter gourd, okra, a drumstick, maize and chickpeas.}

Complete the last lesson of chapter 259: I can say a bitter gourd, okra, a drumstick, maize and chickpeas.
\end{chapteropening}

\section[\texorpdfstring{\ml{പാവയ്ക്ക} (pāvaykka)}{പാവയ്ക്ക (pāvaykka)}]{\ml{പാവയ്ക്ക} (pāvaykka) — a bitter gourd}

\label{lesson:ML-C259-pavaykka}



\begin{quote}
Before the new one: say the Malayalam for a square, then the Malayalam for a triangle.
\end{quote}

\subsection*{You'll want to know: \ml{പാവയ്ക്ക}}
\textbf{\ml{പാവയ്ക്ക}} — \emph{pāvaykka} — \textquotedblleft{}a bitter gourd\textquotedblright{}.

\begin{grammarlens}[title={What you've built}]
One more word for this chapter.
\end{grammarlens}

\subsection*{Your turn}
\begin{itemize}
\raggedright
  \item \emph{Say it:} \emph{pāvaykka}
  \item \emph{Say it:} \emph{pāvaykka}, once more
  \item \emph{Say it:} say \emph{pāvaykka}
  \item \emph{Recall:} say the Malayalam for a square, then the Malayalam for a triangle, then say \emph{pāvaykka} again
\end{itemize}

\begin{tcolorbox}[breakable,colback=teal!4,colframe=teal!35!black,title={Before you move on}]
What does \ml{പാവയ്ക്ക} mean? (\textquotedblleft{}A bitter gourd\textquotedblright{}.) Say it once more.
\end{tcolorbox}
\section[\texorpdfstring{\ml{വെണ്ടയ്ക്ക} (veṇṭaykka)}{വെണ്ടയ്ക്ക (veṇṭaykka)}]{\ml{വെണ്ടയ്ക്ക} (veṇṭaykka) — okra, lady's finger}

\label{lesson:ML-C259-ventaykka}



\begin{quote}
Before the new one: say the Malayalam for a triangle, then the Malayalam for a bitter gourd.
\end{quote}

\subsection*{You'll want to know: \ml{വെണ്ടയ്ക്ക}}
\textbf{\ml{വെണ്ടയ്ക്ക}} — \emph{veṇṭaykka} — \textquotedblleft{}okra, lady's finger\textquotedblright{}.

\begin{grammarlens}[title={What you've built}]
One more word for this chapter.
\end{grammarlens}

\subsection*{Your turn}
\begin{itemize}
\raggedright
  \item \emph{Say it:} \emph{veṇṭaykka}
  \item \emph{Say it:} \emph{veṇṭaykka}, once more
  \item \emph{Say it:} say \emph{veṇṭaykka}
  \item \emph{Recall:} say the Malayalam for a triangle, then the Malayalam for a bitter gourd, then say \emph{veṇṭaykka} again
\end{itemize}

\begin{tcolorbox}[breakable,colback=teal!4,colframe=teal!35!black,title={Before you move on}]
What does \ml{വെണ്ടയ്ക്ക} mean? (\textquotedblleft{}Okra, lady's finger\textquotedblright{}.) Say it once more.
\end{tcolorbox}
\section[\texorpdfstring{\ml{മുരിങ്ങക്ക} (muriṅṅakka)}{മുരിങ്ങക്ക (muriṅṅakka)}]{\ml{മുരിങ്ങക്ക} (muriṅṅakka) — a drumstick (the vegetable)}

\label{lesson:ML-C259-murinnakka}



\begin{quote}
Before the new one: say the Malayalam for a bitter gourd, then the Malayalam for okra.
\end{quote}

\subsection*{You'll want to know: \ml{മുരിങ്ങക്ക}}
\textbf{\ml{മുരിങ്ങക്ക}} — \emph{muriṅṅakka} — \textquotedblleft{}a drumstick (the vegetable)\textquotedblright{}.

\begin{grammarlens}[title={What you've built}]
One more word for this chapter.
\end{grammarlens}

\subsection*{Your turn}
\begin{itemize}
\raggedright
  \item \emph{Say it:} \emph{muriṅṅakka}
  \item \emph{Say it:} \emph{muriṅṅakka}, once more
  \item \emph{Say it:} say \emph{muriṅṅakka}
  \item \emph{Recall:} say the Malayalam for a bitter gourd, then the Malayalam for okra, then say \emph{muriṅṅakka} again
\end{itemize}

\begin{tcolorbox}[breakable,colback=teal!4,colframe=teal!35!black,title={Before you move on}]
What does \ml{മുരിങ്ങക്ക} mean? (\textquotedblleft{}A drumstick (the vegetable)\textquotedblright{}.) Say it once more.
\end{tcolorbox}
\section[\texorpdfstring{\ml{ചോളം} (cōḷaṁ)}{ചോളം (cōḷaṁ)}]{\ml{ചോളം} (cōḷaṁ) — maize, corn}

\label{lesson:ML-C259-colam}



\begin{quote}
Before the new one: say the Malayalam for okra, then the Malayalam for a drumstick.
\end{quote}

\subsection*{You'll want to know: \ml{ചോളം}}
\textbf{\ml{ചോളം}} — \emph{cōḷaṁ} — \textquotedblleft{}maize, corn\textquotedblright{}.

\begin{grammarlens}[title={What you've built}]
One more word for this chapter.
\end{grammarlens}

\subsection*{Your turn}
\begin{itemize}
\raggedright
  \item \emph{Say it:} \emph{cōḷaṁ}
  \item \emph{Say it:} \emph{cōḷaṁ}, once more
  \item \emph{Say it:} say \emph{cōḷaṁ}
  \item \emph{Recall:} say the Malayalam for okra, then the Malayalam for a drumstick, then say \emph{cōḷaṁ} again
\end{itemize}

\begin{tcolorbox}[breakable,colback=teal!4,colframe=teal!35!black,title={Before you move on}]
What does \ml{ചോളം} mean? (\textquotedblleft{}Maize, corn\textquotedblright{}.) Say it once more.
\end{tcolorbox}
\section[\texorpdfstring{\ml{കടല} (kaṭala)}{കടല (kaṭala)}]{\ml{കടല} (kaṭala) — chickpeas}

\label{lesson:ML-C259-katala}



\begin{quote}
Before the new one: say the Malayalam for a drumstick, then the Malayalam for maize.
\end{quote}

\subsection*{You'll want to know: \ml{കടല}}
\textbf{\ml{കടല}} — \emph{kaṭala} — \textquotedblleft{}chickpeas\textquotedblright{}.

\begin{grammarlens}[title={What you've built}]
One more word for this chapter.
\end{grammarlens}

\subsection*{Your turn}
\begin{itemize}
\raggedright
  \item \emph{Say it:} \emph{kaṭala}
  \item \emph{Say it:} \emph{kaṭala}, once more
  \item \emph{Say it:} say \emph{kaṭala}
  \item \emph{Recall:} say the Malayalam for a drumstick, then the Malayalam for maize, then say \emph{kaṭala} again
\end{itemize}

\begin{tcolorbox}[breakable,colback=teal!4,colframe=teal!35!black,title={Before you move on}]
What does \ml{കടല} mean? (\textquotedblleft{}Chickpeas\textquotedblright{}.) Say it once more.
\end{tcolorbox}
